\chapter{AI Augmentation Pipeline}\label{ch:ai}

\section{Architecture}

The AI pipeline is a five-stage request-response chain, hosted on the
internal vLLM cluster per the \texttt{vllm-only} routing policy in
\path{docs/tb/structured/data-products/trial-balance-review-v1.json}.
Every stage reads and writes records validated by \cref{ch:schema};
no blob passes between stages without a schema binding.

\begin{figure}[htbp]
\centering
\begin{tikzpicture}[
  node distance=6mm and 4mm,
  stage/.style={draw,rounded corners=2pt,minimum height=10mm,
                minimum width=22mm,align=center,font=\scriptsize,
                fill=tbblue!15,draw=tbblue!60},
  io/.style={draw,minimum height=7mm,minimum width=20mm,
             align=center,font=\tiny\ttfamily,
             fill=tbgreen!12,draw=tbgreen!50},
  arr/.style={-{Latex[length=2mm]},thick}
]
  \node[stage] (ingest)  {1. Ingest \\ \scriptsize PSGL / S/4};
  \node[stage,right=of ingest]  (anomaly) {2. Anomaly \\ detection};
  \node[stage,right=of anomaly] (commentary) {3. Commentary \\ drafting};
  \node[stage,right=of commentary] (validate) {4. Validate \\ and route};
  \node[stage,right=of validate]  (narrative) {5. FORTM \\ narrative};

  \node[io,below=of ingest]     (rec1) {TB extract};
  \node[io,below=of anomaly]    (rec2) {Variance \\ record};
  \node[io,below=of commentary] (rec3) {Commentary \\ draft};
  \node[io,below=of validate]   (rec4) {Routed \\ summary};
  \node[io,below=of narrative]  (rec5) {FORTM \\ section};

  \foreach \a/\b in {ingest/anomaly,anomaly/commentary,
                     commentary/validate,validate/narrative}
     \draw[arr] (\a) -- (\b);
  \foreach \s/\r in {ingest/rec1,anomaly/rec2,commentary/rec3,
                     validate/rec4,narrative/rec5}
     \draw[arr] (\s) -- (\r);
\end{tikzpicture}
\caption{Five-stage AI pipeline; each stage emits a schema-bound
artefact.}\label{fig:ai-pipeline}
\end{figure}

\section{Stage 1: Ingestion}

\begin{description}
  \item[Input.] PSGL master account data and the SC~Bridge TB
        extract, one per LE, keyed on the \texttt{UPDATE\_TB\_DETAILS}
        workflow state.
  \item[Transform.] Column normalisation against
        \path{docs/tb/machine-readable/training-data/hkg-tb-schema.yaml}
        (13 sheets). Bank CSV samples live in
        \path{docs/tb/extracted-text/hkg-tb-sample.csv} and
        \path{docs/tb/extracted-text/hkg-pl-sample.csv}.
  \item[Acceptance.] \texttt{EUC-001} evidence rules apply: the
        extraction must preserve debit/credit totals and every row
        must carry a valid account number from the refreshed master.
\end{description}

\section{Stage 2: Anomaly detection}

The anomaly detector performs four orthogonal checks per account
(\cref{tab:anomalies}). Thresholds are account-category-sensitive and
are re-fit monthly from the prior twelve months; the fitted thresholds
are persisted to the variance file EUC as audit evidence.

\begin{table}[htbp]
\caption{Anomaly checks run per account.}\label{tab:anomalies}
\centering
\small
\begin{tabularx}{\linewidth}{@{}llX@{}}
\toprule
\textbf{Check} & \textbf{Metric} & \textbf{Rule} \\ \midrule
Period spike   & $|\Delta|$       & $> 3\sigma$ over prior 12 months$^{\dagger}$ \\
Sign flip       & sign($v_t$)       & $\neq$ sign($v_{t-1}$) and $|v_t| > $ threshold \\
New account     & existence         & account absent from prior-period master \\
Stale balance   & $|\Delta|$        & $= 0$ for $\ge 3$ consecutive periods with prior activity \\
\bottomrule
\end{tabularx}
\end{table}

\noindent$^{\dagger}$\textbf{Assumption caveat:} The $3\sigma$ threshold assumes
normally-distributed variance values. This assumption has not been statistically
validated against the TB population; financial data often exhibits fat tails
(leptokurtosis). The pilot will include a Shapiro-Wilk normality test per account
category. Non-normal distributions may require robust alternatives (e.g., MAD-based
z-scores or quantile thresholds). Thresholds marked ``TBD pending pilot.''

Only accounts flagged by at least one check enter Stage 3. Accounts
below materiality (\texttt{DP-001}) are still logged for audit but
short-circuit to Stage 4.

\section{Stage 3: Commentary drafting}

\paragraph{Model.} A vLLM-hosted instruction-tuned model accessed
via the internal gateway. The system prompt is the
\texttt{x-prompting-policy.systemPrompt} field of the data product
record; \texttt{maxTokens = 4096}, \texttt{temperature = 0.3},
\texttt{responseFormat = "text"}.

\paragraph{Prompt structure.} Each commentary request carries:
\begin{itemize}
  \item The flagged variance record (account, prior, current, delta).
  \item The three most similar historical commentaries retrieved from
        the \texttt{TB\_REVIEW\_VECTORS} table (55 chunks, 220-word
        chunking with 40-word overlap; see
        \corpus\texttt{machine-readable/rag/rag\_embedding\_records.jsonl}).
  \item The account's finance caption and segment.
  \item Any related journal identifiers found in the GL for the period.
\end{itemize}

\paragraph{Output contract.} The draft is returned as a
\texttt{commentary\_draft} blob that is paired with the variance
record before leaving the stage. Human review by the GFS Analyst is
mandatory (per DOI \texttt{REVIEW\_INCORPORATE}); the draft is a
\emph{suggestion}, never a final artefact.

\section{Stage 4: Validation and routing}

The gateway logic defined in \cref{ch:workflow} is evaluated here.
\Cref{tab:routing} maps each decision point to the record fields the
engine reads.

\begin{table}[htbp]
\caption{Decision-point routing.}\label{tab:routing}
\centering
\small
\begin{tabularx}{\linewidth}{@{}llX@{}}
\toprule
\textbf{DP} & \textbf{Gateway state} & \textbf{Read fields} \\ \midrule
DP-001 & CHECK\_MATERIAL\_VARIANCES & variance\_amount, account\_category, materiality\_limit \\
DP-002 & IS\_JOURNAL\_POSTED         & journal\_status, posting\_date \\
DP-003 & CHECK\_FURTHER\_INVESTIGATION & investigation\_status, entry\_age\_days \\
DP-004 & CHECK\_CONFIRMATION         & confirmation\_received, days\_to\_deadline (M+4) \\
DP-005 & CHECK\_UNEXPLAINED          & unexplained\_variance\_count \\
\bottomrule
\end{tabularx}
\end{table}

\section{Stage 5: FORTM narrative generation}

The final stage composes the FORTM pack section. It consumes the
routed summary plus the KPI telemetry (\cref{tab:kpis}) and emits a
structured narrative containing: executive summary, material variance
table, unexplained items, risk register update, and the list of
outstanding journal postings. The narrative goes to
\stateid{REVIEW\_INCORPORATE} for \texttt{EUC-003} sign-off before
\stateid{SEND\_FORTM}.

\section{Controls posture of the pipeline}

Every stage is in-scope of the UK ACG automated-controls regime. The
per-stage control posture is summarised in \cref{tab:ai-controls}.

\begin{table}[htbp]
\caption{Per-stage controls posture.}\label{tab:ai-controls}
\centering
\small
\begin{tabularx}{\linewidth}{@{}llX@{}}
\toprule
\textbf{Stage} & \textbf{Control type} & \textbf{Mechanism} \\ \midrule
Ingest        & preventive & Totals reconciliation vs PSGL source \\
Anomaly       & detective  & Threshold log + human sampling \\
Commentary    & detective  & Mandatory human review (DOI) \\
Validation    & preventive & Schema + gateway enforcement \\
FORTM         & detective  & EUC-003 senior sign-off \\
\bottomrule
\end{tabularx}
\end{table}

\section{Evaluation harness}

\begin{description}
  \item[Golden set.] The HK Nov~2025 workbooks
        (\corpus\texttt{source/HKG\_TB review Nov'25.xlsm} and
        \corpus\texttt{source/HKG\_PL review Nov'25.xlsm}) serve as
        regression data. All 234 LEs enter at phase~4 of the
        implementation plan.
  \item[Commentary metric.] Draft quality is scored against analyst
        accepts in the DOI review phase; target accept-without-edit
        $\ge$ 70\%, accept-with-edit $\ge$ 95\%.
  \item[Anomaly metric.] Per-check precision/recall against analyst
        flags in a held-out period.
  \item[Cycle time.] Measured as \texttt{TB-KPI-002}; target is the
        enablement of a daily cadence, which collapses the cycle from
        \texttt{M+1..M+7} to a rolling one-day delay.
\end{description}
